\section{The Coupling Framework}
\label{sec:framework}

The interface transfers a small set of model outputs and policy inputs into explicit
changes in the receiving model. Each named mapping records the source series, its economic
interpretation, the receiving-model entry, the applicable accounting rule, and the tests
that have been completed. This section describes those common elements, the deployed
electricity-cost signal (\Cref{sec:signal}), its transmission into an economy without an
intermediate energy input (\Cref{sec:routes}), the staged architecture and country
contract (\Cref{sec:architecture}), and the accounting rules that prevent overlapping
representations (\Cref{sec:discipline}).

\subsection{Transferred quantities and receiving-model entries}
\label{sec:signals-wedges}

The two models meet across an asymmetry. \OGCore{} is rich in behaviour, public finance,
and household heterogeneity but blind to energy technology; \CLEWS{} is rich in
technology, dispatch, and emissions but blind to behaviour, the fiscal sector, and
distribution. From the resource model come an electricity-cost path, investment
requirements, and emissions. From the economy come sectoral activity, final demand, and
an equilibrium return. Policy scenarios also supply a carbon-price path or investment
incentive to the relevant model entries. The receiving entries include the
consumption-tax vector $\tauc$, industry productivities $\TFP_m$, the public-investment
share $\alphaI$, demographic objects $\rho(s,t)$ and $e(t,s,j)$, and \CLEWS{} demand and
discount-rate parameters. Some are literal wedges; others are ordinary parameters or
model inputs. The interface checks dimensions, admissible ranges, units, and
concordances before applying them. Portability is conditional on the country contract
identifying the corresponding industries, goods, commodities, and units
(\Cref{sec:architecture}).

\subsection{The electricity-cost signal}
\label{sec:signal}

The theoretically preferred electricity signal is a marginal cost---the \emph{shadow
price} on the \OSeMOSYS{} commodity-balance
constraint \citep{howells_osemosys_2011} (which requires supply to meet demand in every
timeslice and year), what the least-cost system would pay to deliver one more unit, and
under the competitive interpretation the price that clears the market. But that shadow
price is degenerate in this model: it binds in only scattered years and is near-zero
elsewhere, so it is unusable as a smooth long-run signal (it is kept for scarcity
studies). The coupling therefore uses the reform/base ratio of a curated cost index or,
where that index is unavailable, reconstructed \emph{levelized cost of electricity}
(LCOE): the average capital, operation and maintenance, and fuel cost per unit of
generation. This is a dense average-cost proxy for the price the economy faces, not a
marginal or observed retail price. The proxy performs two functions.

First, it closes an elasticity gap. Inside the energy optimiser, demand is an
exogenous datum.

\begin{assumption}[Price-inelastic energy demand in \CLEWS{}]
\label{ass:inelastic}
\CLEWS{}/\OSeMOSYS{} solves for the least-cost system that meets a \emph{fixed} energy-
service demand; quantity demanded does not respond to the commodity price within the
optimiser.
\end{assumption}

\noindent The behavioural response must therefore come from the macro side. Feeding
this cost proxy to \OGCore{} households---whose multi-good demand is unit-elastic above a
Stone--Geary subsistence floor \citep{stone_1954}, see \Cref{sec:og}---makes the
quantity of energy demanded respond to its price, supplying the elasticity the
optimiser lacks.

Second, it reveals incidence. A single program-level ``cost'' is silent on who bears
it. \OGCore{}'s households, indexed by $\Ngen$ age and $\Jinc$ lifetime-income groups,
distribute the burden of an energy-cost change. With the Stone--Geary calibration used
here, lower-income households devote a larger budget share to the energy good, permitting
the model to test whether incidence is regressive; a representative household cannot
report that distribution.

These two roles define what it means for the coupling to be internally consistent.

\begin{definition}[Soft-link equilibrium]
\label{def:fixedpoint}
Let $D_{\text{OG}}(\pi)$ be the energy demand \OGCore{} implies when households (and,
in the structural form of \Cref{sec:routes}, firms) face energy price $\pi$, and let
$\pi_{\text{C}}(D)$ be the \CLEWS{} price of delivering $D$ (its levelized cost). A \emph{consistent} coupling is a pair $(\pi^\star, D^\star)$ with
\[
  D^\star = D_{\text{OG}}(\pi^\star)
  \qquad\text{and}\qquad
  \pi^\star = \pi_{\text{C}}(D^\star).
\]
\end{definition}

\noindent A single \CLEWS{}$\to$\OGCore{} pass evaluates $\pi_{\text{C}}(D^{b})$ at
\emph{baseline} demand and applies it once; iterating the two maps to the fixed point
of \Cref{def:fixedpoint} is \emph{loop closure} (\Cref{sec:architecture}). With the
levelized cost as the deployed price, the fixed point is consistency under that
pricing proxy, not a competitive-equilibrium price---the distinction inherits from
the signal choice above. The
present work tests the single pass and, separately, the \CLEWS{} re-solve seam that the
closed loop requires; the emitted paths have not yet been carried through that seam by an
outer iteration controller.

\begin{remark}[Shadow-price hygiene]
\label{rem:hygiene}
The failed dual motivates a protocol applied to any marginal cost read from the optimiser
before it crosses the seam: the value must sit above the solver's own reporting
resolution; it must be invariant across alternate optima---a price that moves while
nothing physical changes is degeneracy, not information; it must be decomposed
against any cost parameter that floors it; and it must trace to a sourced input.
The electricity marginal cost fails the first two tests, which motivates the
levelized-cost proxy above. Preliminary work on the Philippine land balance finds the
same failure modes; that application is not part of the evidence reported here.
\end{remark}

\subsection{Transmission of electricity costs}
\label{sec:routes}

The design follows from identifying the purchasers exposed to the electricity-cost
change. In the Philippine social accounting matrix they have two parts: industries,
which buy roughly three quarters of electricity as a production input, and households,
which consume the remaining quarter directly. A faithful transmission must reach both
groups of buyers while applying the increase to disjoint expenditure flows.

The difficulty is that both buyers exist in \OGCore{}, but intermediate purchases do
not. Two structural facts matter. First, production is value added only: each industry
produces with private capital, public capital, and labour (\Cref{sec:og}), and
\emph{no firm buys anything from any other firm}---there is no line in the
manufacturing sector's cost function where electricity appears, at any price. Second,
the model's input--output object is a \emph{consumption} bridge, not an
intermediate-use matrix: a fixed $\Igood\times\Mind$ map that converts industry
outputs into consumption goods, $p_i=\sum_m \omega_{im}p_m$, with no substitution
across industries. Electricity's price can therefore reach households through the
goods bridge, but the three quarters of the increase that industrial buyers should
pay has no structural path---the transactions are missing from the model. The two-route
approximation below follows from that omission.

\paragraph{A producer-price route does not reach intermediate users.}
One possible representation changes the electricity
\emph{producer}: let the industry price $p_{m_e}$ itself rise, and let households
respond to that. \OGCore{} has no parameter that inserts a wedge between what a
firm's output sells for and what the firm receives---prices are outcomes, not
dials---but productivity provides a reduced-form handle. Because every industry prices at unit
cost (zero profit), $p_m$ is proportional to $1/\TFP_m$ holding factor prices, so
dividing the electricity industry's productivity by the cost-proxy ratio,
$\TFP_{m_e}\to\TFP_{m_e}/g_t$, makes the model \emph{produce} the higher electricity
price endogenously---with the full general-equilibrium response of wages, the
interest rate, and the sector's own inputs that a tax wedge does not have. This route
is implemented as a controlled comparison. Its transmission is limited.
It flows through the consumption bridge into the energy good, and households pay it.
It goes nowhere else: manufacturing's costs, prices, and decisions are untouched,
because manufacturing buys no electricity in the model. The route is not wrong---it is
\emph{incomplete by construction}: it can deliver the shock only to household final
use, roughly one quarter of Philippine electricity expenditure. It also
encodes a \emph{system}-cost change as an electricity-\emph{technology} change, so the
sector's own capital and labour reorganise in response to a productivity event that
did not occur in the scenario.

\paragraph{Household final use: a revenue-neutral consumption wedge.}
The deployed channel carries household final use through a consumption-tax wedge on the
energy good $i_e$, set so the gross-of-tax price inherits the cost-proxy ratio,
\begin{equation}
  1 + \tau^{c,r}_{t,i_e} \;=\; g_t\,\bigl(1 + \tau^{c,b}_{t,i_e}\bigr),
  \label{eq:routeA}
\end{equation}
where $g_t \equiv \pi^{r}_t/\pi^{b}_t$ is the \CLEWS{} cost-proxy ratio. Households then
face $(1+\tauc_{t,i_e})\,p_{i_e,t}$ exactly where \OGCore{} puts consumption taxes:
in the cost-of-living index, in subsistence spending, and in the sub-good demand.
Demand falls and, under the Stone--Geary calibration used here, the incidence is
regressive. The wedge raises what households pay without representing a technology event
inside the electricity sector. A consumption tax also credits the government with
revenue, whereas the electricity-cost increase represented here is not government
revenue. Left alone, the wedge would therefore introduce a fiscal consolidation in
addition to the price change. The channel neutralises that effect by returning the revenue to
households as a lump-sum transfer (a bump to the transfer share $\alphaT$, estimated
on baseline quantities): the wedge is \emph{revenue-neutral by construction},
delivering relative prices and incidence and nothing else. This side alone is close
to output-neutral: it reshuffles consumption rather than shrinking output
(\Cref{sec:results}).

\paragraph{Intermediate use: a reduced-form industry cost-push.}
For intermediate purchases, the channel raises each industry's costs directly by the
amount implied by its electricity exposure. Let
$\boldsymbol\theta=(\theta_1,\dots,\theta_{\Mind})$ be the direct electricity cost
shares and $A$ the matrix of interindustry input coefficients, both read from the
social accounting matrix. A proportional price increase $\hat p_E$ raises industry
$m$'s unit cost directly by $\theta_m\hat p_E$ and, through purchased intermediates,
indirectly; summing the Leontief series gives the total cost-push
\begin{equation}
  \hat{\mathbf c} \;=\; (I - A^{\!\top})^{-1}\,\boldsymbol\theta\,\hat p_E ,
  \label{eq:routeB}
\end{equation}
applied as a productivity haircut $\TFP^{r}_m = \TFP^{b}_m\,(1-\kappa_m)$, with
$\kappa_m$ the gross-output cost increase of industry $m$ and electricity's own
self-use netted out so that nothing appears in both legs. The \emph{size} of the
shock comes from data; the transmission mechanism is imposed because intermediate inputs
are absent from the production block. It approximates the first-order cost pass-through
of an intermediate-input structure: firms do not substitute away from electricity, and
no intermediate energy market clears. The productivity reduction represents the resource
cost in reduced form; it is not a literal loss of productive capacity.

\begin{remark}[The two routes cover disjoint expenditure bases]
\label{rem:ab-contrast}
Households and firms face the same proportional electricity-cost change; the social
accounting matrix determines the expenditure base exposed through each route. The
consumption wedge carries final-use prices and incidence and is fiscally neutralised by
the rebate. The cost-push carries intermediate exposure through
$\boldsymbol\theta$. Electricity self-use is removed before upstream propagation, so
the two expenditure bases do not overlap. Neither route alone represents the complete
electricity-cost mapping.
\end{remark}

\paragraph{Structural extension: energy as an intermediate input.}
A structural alternative would add firms' electricity purchases by nesting energy with
value added in the technology,
\begin{equation}
  Y^{\text{gross}}_m
  \;=\; \TFP_m\Bigl[\,\omega_m\, VA_m^{\frac{\sigma_{E,m}-1}{\sigma_{E,m}}}
        + (1-\omega_m)\, E_m^{\frac{\sigma_{E,m}-1}{\sigma_{E,m}}}\Bigr]
        ^{\frac{\sigma_{E,m}}{\sigma_{E,m}-1}},
  \label{eq:routeC}
\end{equation}
where $VA_m$ is the existing capital--labour aggregate and $E_m$ is energy purchased
from the energy industry $m_e$ at the price $p_{m_e}$. Cost minimisation yields
the energy demand
\begin{equation}
  E_m \;=\; \theta_m\Bigl(\frac{p_{m_e}}{MC_m}\Bigr)^{-\sigma_{E,m}} Y^{\text{gross}}_m ,
  \label{eq:routeC-foc}
\end{equation}
the energy industry clears a resource constraint
$Y_{m_e} = (\text{final demand}) + \sum_{m\ne m_e} E_m$, output is measured as value
added net of energy purchases, and the energy price feeds firm cost endogenously so
that the fixed point of \Cref{def:fixedpoint} closes the loop on quantities as well as
prices. The substitution elasticities $\sigma_{E,m}\in[0.3,0.8]$ are available from the
energy-CGE literature and the cost shares $\theta_m$ are already calibrated from the
Philippine social accounting matrix (\Cref{sec:calibration}); indeed the SAM already
contains the full intermediate-use block, so the barrier is \OGCore{}'s firm side,
not data. With it, the two-legged composite collapses into one mechanism---firms
substitute away from expensive electricity, an energy market clears, and the fixed
point of \Cref{def:fixedpoint} closes on quantities as well as prices. It is a
production-side extension to \OGCore{} itself, out of scope here.

\subsection{Architecture: a staged soft link}
\label{sec:architecture}

The classic joint-planner hard link is unavailable without reformulating both models.

\begin{assumption}[Incompatible solution concepts]
\label{ass:hardlink}
\OGCore{} solves an overlapping-generations transition path period by period under
foresight; \OSeMOSYS{} is a single perfect-foresight linear program over all years at
once. \OGCore{} exposes no scalar welfare into which the program could be folded, and
no shock can be injected mid-solve.
\end{assumption}

\noindent The assumption forecloses the integrated-welfare-program route of the
hard-link tradition. A simultaneous formulation that stacks both models' equilibrium
and optimality conditions as a complementarity system is conceivable, but it would be a
new model rather than a direct composition of the existing programs.

\noindent The models therefore stay separate and exchange a few signals---a soft link,
in the tradition of iteratively coupled energy--macro models
\citep{messner_schrattenholzer_2000,manne_merge_1995,loulou_times_2008}, but with a
behaviourally and fiscally rich overlapping-generations core in place of a
representative agent. The orchestration (\Cref{fig:architecture}) runs in stages:
(i)~solve the \OGCore{} baseline in steady state and along the transition;
(ii)~solve the \CLEWS{} scenario and extract the forward quantities; (iii)~apply the
receiving-model changes and solve the reform transition; (iv)~read the
\OGCore{}$\to$\CLEWS{} quantities (the equilibrium return $\rp$ and sectoral activity);
and, in the closed loop, (v)~re-solve \CLEWS{} and repeat to the fixed point of
\Cref{def:fixedpoint}.

\diagramfig{architecture}{The structural seam. The two models stay separate and pass
quantities across a contract; the classic joint-planner formulation is unavailable under
\Cref{ass:hardlink}.}{fig:architecture}

What makes the exchange auditable is an explicit interface contract:
a \emph{scenario pair} naming the baseline and reform runs, the
\CLEWS{} scenario, and the country; a \emph{concordance} mapping \OGCore{} industries
$m$ and goods $i$ to \CLEWS{} commodities, technologies, and demands---supplied by the
country configuration, since \OGCore{} has no native notion of a resource; and a
\emph{unit map} (real versus nominal, base year, deflator) that places \CLEWS{} quantities
in the \OGCore{} real numéraire. The interface can be re-pointed only where this contract
identifies compatible country-specific objects.

Loop closure is constrained by \Cref{ass:inelastic}. Because energy demand is
price-inelastic inside the optimiser, naïve price feedback can oscillate; convergence
to \Cref{def:fixedpoint} requires damping and a baseline-anchored demand calibration.
The \OGCore{}$\to$\CLEWS{} emitters produce well-formed \CLEWS{} inputs, and the
\CLEWS{} re-solve seam they are designed to feed---step~(v) of the orchestration---is
implemented and tested with a controlled perturbation: the link
patches a copy of the \CLEWS{} case, re-solves it
in the energy model's own environment, and reads the result, leaving the calibrated case
untouched. On the Philippine case a controlled $+10\%$ household-electricity demand
patch propagates through the re-solve to a $+4.3\%$ electricity-production response,
confirming that the seam transmits a demand change. The emitted return and activity paths
have not been fed through it end to end. What is not yet wired is the
outer controller that iterates the re-solve to \Cref{def:fixedpoint} under damping; until
it runs, the two \OGCore{}$\to$\CLEWS{} channels (\Cref{sec:ch-discount,sec:ch-demand})
emit well-formed inputs but have no feedback effect on the macro economy.

\subsection{No-double-counting discipline}
\label{sec:discipline}

Because the same source quantity can be represented through more than one parameter---and because a
fiscal instrument used as a price wedge also moves the budget---the framework needs
explicit accounting rules. We state them as principles; each is enforced by a guard in
the corresponding channel where the forbidden combination can arise within one channel,
and by the composition rules of the experiment layer where it arises across channels.

\begin{principle}[One representation of each carbon cost]
A carbon penalty that changes the \CLEWS{} solution may already enter \OGCore{} through
the transmitted electricity-cost ratio. The same electricity cost is not added again as
a separate macroeconomic wedge. A separate \OGCore{} carbon wedge is used only for a
carbon-bearing household expenditure not already represented in that ratio, and an
implied shadow price is not stacked on an exogenous carbon-price path.
\end{principle}

\begin{principle}[Count each energy cost once]
Represent each unit of an energy cost once. The energy channel (\Cref{sec:routes})
partitions electricity by use---a consumption wedge on the household share, a cost-push
on the intermediate share, with electricity's own self-use netted out---so the two sides
do not overlap. What is forbidden is a second representation of the \emph{same} flow: a
$\tauc$ wedge and a $\TFP$ haircut on the same consumption, or a separately-inferred
carbon wedge stacked on top.
\end{principle}

\begin{principle}[Generation-capital representations require distinct sources]
Multiple parameters are not calibrated to the same observed generation-capex change. A
technology-driven change in the capital share and a separately specified investment tax
credit may coexist when the scenario identifies both effects; they are not alternative
names for one source quantity. In the Cobb--Douglas core the capital-share lever changes
factor shares, whereas the ITC changes the user cost of capital
(\Cref{sec:ch-capint}).
\end{principle}

\begin{principle}[Revenue must be routed]
A $\tauc$ wedge mechanically generates consumption-tax revenue. That revenue must be
recycled (as transfers or public investment) or explicitly financed (as deficit), or
the run introduces an additional fiscal change alongside the price shock.
\end{principle}

\begin{principle}[Public and private capital are distinct]
Public-grid capital (reaching \OGCore{} through the public-investment share $\alphaI$
and the public-capital stock $\Kg$) and private-generation capital (reaching the
firm's own capital) touch different objects. Coupling both is complementary, not double
counting.
\end{principle}
